% Source: 41_Chapter_31_Unwilling_to_Die_for_an_Old_Idea_Denomi.txt
\bookchapter[Unwilling to Die for an Old Idea]{CHAPTER 31}{UNWILLING TO DIE FOR AN OLD IDEA}{Denominations}

\leadin{F}{OR DECADES CRITICS HAVE} called them “a scandal,” “a blight,” “factionalism,” and “a
caste system,” but denominations remain the institutional hallmark of modern Christianity.

Criticism is understandable. Any Christian reading his New Testament senses the difference between
the faith of the apostles and the Christianity of our day. The apostle Paul, for example, speaks of
the church as the temple of God, unified in devotion to the Lord Jesus Christ, but we find in our
time only a menagerie of cults, sects, denominations, and isms. We sense deeply that this divided
state of Christianity ought not be, but it is.

Why? How did denominations come to be the primary expression of Christianity in modern times?

The simple fact is Christians are divided today, in part at least, because they have the freedom to
differ. In earlier centuries they did not. We may curse denominations or try to ignore them, but
they are not going to disappear soon because the cost of their removal is greater than most of us
want to pay. We are shocked by this fruit of modern Christianity, but few of us want to lay the ax
at its root.

The Age of the Reformation (1517–1648) did not suddenly end and the Age of Reason and Revivals
(1648–1789) appear overnight as historians’ dates might suggest, but times do change, and one marked
difference between the sixteenth century and the seventeenth was the acceptance of religious
differences.

We often hear, “I disapprove of what you say, but I will defend to the death your right to say it.”
And most Christians today, even if they cannot identify its source, accept the sentiment of that
oft-quoted manifesto. They accept it not because it is Christian but because it is modern.

\section{SUPPRESSION OF THE NONCONFORMIST}

The statement probably belongs to Voltaire (1694–1778), the proud, self-sufficient humanist of the
Age of Reason. It is the kind of thing Martin Luther or Ignatius Loyola would never say, because
neither believed it. In the light of the Reformation, dissent was neither a Christian virtue nor a
human right. The Reformers were as eager as Catholics to suppress nonconformity. That was because
both camps believed that Christian truth held societies together. It was an instrument of power. And
only one side in a religious conflict had the truth. The idea that God’s Word could be found on both
sides of a battle line was a revolutionary concept that only gained a hearing after both sides fell
from exhaustion.

In the 1540s and 1550s, when both Protestants and Catholics still believed in a united Christendom,
Lutheran princes fought Catholic Imperial forces to a stalemate in Germany. In the Peace of Augsburg
(1555) both sides agreed to stop fighting only after adopting the territorial principle: the ruler
in each region could decide the faith of his subjects, Lutheran or Catholic.

The compromise was a spin-off of the idea of the Christian prince. It was a guarantee to a number of
people the right to worship according to their consciences, but it also recognized the right of the
prince to persecute those who didn’t happen to agree with him. It led to suffering and hardship for
many innocent people whose only offense was religious dissent.

The principle of territorialism was the herald of the approaching death of Christendom. If one
Christian prince of a small territory can lay down the law of one religion for his subjects and
another ruler a dozen miles away can lay down another law for another religion, then the norm of
Christian truth is shattered and Christian society crumbles to pieces.

Of course, events were running ahead of new ideas, so other wars were fought all over Europe and new
territories were fashioned according to the old concept that only one religion could exist in one
region.

From 1562 to 1598 France suffered a series of civil wars between Roman Catholics and French
Calvinists (or Huguenots). When both parties reached the point of utter fatigue they agreed to a
territorial compromise in the royal Edict of Nantes (1598). The Huguenots gained religious freedom
and political control of certain parts of the country while Roman Catholicism remained the official
religion of the realm and in the greater portion of the nation.

In a similar way between 1560 and 1618 in the Netherlands, the strongly Calvinistic Dutch fought a
war of independence from Catholic Spain and won. In the southern territories, however, the area we
call Belgium, the people remained Catholic and did not gain their independence from Spain until much
later.

\section{THE THIRTY YEARS’ WAR}

All of these conflicts were a bloody prelude to the last and most devastating of the so-called wars
of religion: the Thirty Years’ War (1618–1648). This conflict began primarily as a religious
struggle with political overtones and ended as a barbarous political power struggle with religious
overtones. That shift in dominant motives makes the Thirty Years’ War a fitting symbol of the
transition from the Age of Reformation to the Age of Reason and Revivals.

The war was as complex as it was long. This is no place for military or political details. We are
interested in the change in ideals, and we can find this in the simplest highlights.

One of the glaring weaknesses of the Peace of Augsburg (1555) was that it completely ignored the
Calvinists. Given their sense of holy mission, it was merely a matter of time before hostilities
erupted anew. Preparations for war were laid early in the seventeenth century when Protestants
formed a league of German princes and the Catholics created a similar Catholic League.

Fighting broke out in 1618.

A zealous supporter of the Counter Reformation, the Jesuit-educated Ferdinand II, was named the king
of Bohemia only shortly before he was also elected Holy Roman Emperor. Moving under the concept of
one religion for one realm, Ferdinand attempted to uproot Protestantism from Bohemia and impose
Catholicism upon his subjects.

The Bohemian nobles, mostly Protestants, rose in revolt and offered their crown to Frederick V, the
ardent Calvinist ruler of the Palatinate, one of the major German territories. Frederick’s
acceptance touched off the fighting between Catholics and Calvinists.

In 1620 near Prague the Catholic Imperial forces, surging forward to the cry, “For the Virgin Mary,”
crushed the Bohemians and confiscated most of the estates of the insurgents. The victors added salt
to the wounds by awarding the Jesuits control of the University of Prague, the old haunt of John
Hus.

Appalled by this Catholic victory and eager to annex German territory, King Christian IV of Lutheran
Denmark entered the war against Ferdinand and the Catholic forces. Without adequate support,
however, his venture was doomed from the start. In 1626 the Danes were completely routed in the Harz
Mountains and forced to withdraw to Denmark like a pack of whipped pups.

Imperial control of the southern shores of the Baltic and deep religious convictions compelled the
able Lutheran warrior, King Gustavus Adolphus of Sweden, to enter Germany as the new leader of the
Protestant cause. A series of smashing victories carried him south as far as Munich. “The Lion of
the North” Protestants called him, but even regal courage meets its end. At the Battle of Lutzen
(1632), southwest of Leipzig, the Swedish army was again victorious, but Adolphus was cut down in
combat.

Without Adolphus the war wore on, but the outcome was already clear. The Catholic forces could not
subdue the Protestants in northern Germany and the Protestants could not control the Catholics in
the south.

In the final years of the conflict, religion faded to insignificance. For the most part France and
Spain, both nominally Catholic, struggled for political advantages in the Rhineland.

When the swords fell silent, Germany lay ravaged. Ferdinand’s dream of a revival of imperial
authority there was gone, and in its place were 300 independent states. Out of sheer futility the
religious zeal of Catholics and Calvinists cooled, and men began to question the territorial idea.
Denominationalism was an alternative.

The terms of peace, called the Peace of Westphalia (1648), reflect the passing of an age. Calvinism
joined Lutheranism and Catholicism as a recognized expression of the Christian faith. Princes, if
they chose, could, for the first time, allow Protestants and Catholics to exist within their
territories. And the pope was excluded from any interference in the religious affairs of Germany.
Naturally Pope Innocent X condemned the treaty, but both Catholics and Protestants ignored his
protests. After more than a thousand years the state was free to transact its business as though the
pope did not exist.

Christianity’s alliance with civil power was changing. Christians in the West would now live in
nation-states, the real victor in the Thirty Years’ War. The emerging nation-states would require
Christians to live under a government whose citizens belonged to a variety denominations and faiths.
In time, some of these nation-states required citizens to affirm a generic version of God in order
to hold office or attend college; but distinctive beliefs and practices, such as being Catholic or
Presbyterian, were strictly a private, individual matter.

\section{THE NEW WAY IN AMERICA}

By 1648, far across the Atlantic in the English colonies of North America, men were seeking to build
distinctly religious societies. At first, they came to fashion places of devout Christian character
and witness. Eventually the goal of building a holy city that nurtures holy men was amended to
building a city where people could be holy. Waves of newcomers, the indifference of the second
generation, and the wide open frontier all weighed against a religious uniformity.

During the sixteenth century both England and France had sent courageous explorers to North America
searching in vain for the Northwest Passage to China. But not until the London Company landed its
first colonists at Jamestown, Virginia, in 1607 did permanent settlements begin drawing Englishmen
to the New World. The Pilgrims landed in 1620, and ten years later the Massachusetts Bay Company
began attracting thousands of frustrated Puritans to Boston and its surrounding towns. Between 1629
and 1642 some twenty-five thousand Puritans migrated to New England. Other than in Virginia, the
English authorities made no attempt to impose a pattern of religious uniformity in the New World.

The colonies were commercial ventures designed to contribute to the prosperity of a developing
English empire. To be profitable the colonies needed settlers to clear the forests and till the
fields. So the promise of religious toleration provided a powerful incentive to accept the hazards
of life in the American colonies. Quakers came to Pennsylvania, Catholics to Maryland, and Dutch
Reformed to New York. Later came Swedish Lutherans and French Huguenots, English Baptists and
Scottish Presbyterians.

The prominent exception to this general policy of religious toleration were the Congregational
Puritans of Massachusetts Bay, who were determined to establish a new Zion in the American
wilderness: a “due form of government both civil and ecclesiastical” is the way Governor John
Winthrop put it.

The opportunity to create the rule of the Puritan saints in New England arose from the curious
omission from the charter of the customary clause requiring the headquarters of the company to be in
England and thus subject to the authority of the king. As a consequence of the omission, the
Massachusetts Bay Company, controlled by the saints, was in effect an independent republic, for all
the members of the company migrated to New England and took the charter with them.

Thus for two generations the saints in New England ruled by a policy of religious conformity even as
the same policy proved futile in old England. Individuals could face severe punishment on numerous
grounds: for failing to attend church services, for denying Christ’s resurrection, or for showing
irreverence for the Bible.

Yet even here, within the heart of New England Puritanism, were impulses for religious toleration.
One of them was in the Puritans’ devotion to Scripture.

Pastor John Robinson told the Pilgrim Fathers, “the Lord hath more truth yet to break forth out of
his Holy Word.” The Puritan was ever ready “to receive whatever truth shall be made known~.~.~. from
the written word of God.” Not surprisingly then, New England Puritans were always spawning their own
dissenters, men and women who appealed to the truth that had been made known to them from “the
written word of God.” Roger Williams was one of these; Anne Hutchinson another.

A second check on Puritan intolerance was hidden in the wilderness. Dissenters in New England did
not need to go underground; they could just move on—across the river, through the woods, or over the
mountain. Sanctuary was always possible in open space.

Thus Puritan devotion to the Bible, the presence of the wilderness, and the English policy of
toleration combined to undercut the intolerance of the New England Puritans.

The end of their attempt to enforce conformity in New England was foreshadowed by the loss of their
charter in 1684. The colony’s return to royal control was dramatically signaled in 1687 when the
governor seized the Old South meeting house for Anglican worship. Thereafter New Englanders had to
conform to English policy.

\section{THE IDEA OF DENOMINATIONS}

The religious diversity of the American colonies, though largely within the Puritan tradition,
called for a new understanding of the church. We may call it the denominational theory of the
church. The use of the word denomination to describe a religious group came into vogue about 1740
during the early years of the Evangelical Revival led by John Wesley and George Whitefield. But the
theory itself was hammered out a century before by a group of radical Puritan leaders in England and
America.

Denominationalism, as originally designed, is the opposite of sectarianism. A sect claims the
authority of Christ for itself alone. It believes that it is the true body of Christ; all truth
belongs to it and to no other religion. So by definition a sect is exclusive.

The word denomination by contrast was an inclusive term. It implied that the Christian group called
or denominated by a particular name was but one member of a larger group, the church, to which all
denominations belong.

The denominational theory of the church, then, insists that the true church cannot be identified
with any single ecclesiastical structure. No denomination claims to represent the whole church of
Christ. Each simply constitutes a different form, in worship and organization, of the larger life of
the church.

The Reformers had planted the seeds of the denominational theory of the church when they insisted
that the true church can never be identified in any exclusive sense with a particular institution.
The true succession is not of bishops but of believers. Luther insisted that some institutional
expression “in a place and in the things and activities of the world” was inevitable; but “in this
life the church is not properly understood in terms of all this.” The outward forms of the church,
said Luther, ought to give the Word of God free course in the world, not block its power to save.

In a similar way Calvin, in the preface to his Institutes, indicated that it is impossible to draw
precise boundaries to the church of Christ. No one can determine with precision who is numbered
among the elect of God.

The Reformers, however, never followed this lead. When religious dissent arose within a particular
area, they tried to suppress it. They were still convinced that only one religion could exist in one
region.

The real architects of the denominational theory of the church were the seventeenth-century
Independents (Congregationalists) who represented the minority voice at the Westminster Assembly
(1642–1649). The majority at the Assembly held to Presbyterian principles and expressed these
convictions classically in the Westminster Confession of Faith and in the Westminster Larger and
Shorter Catechisms.

The Independents, however, who held to congregational principles, were keenly aware of the dangers
of “dividing the godly Protestant party” in England, so they looked for some way to express
Christian unity even when Christians did not agree.

These Dissenting Brethren of Westminster articulated the denominational theory of the church in
several fundamental truths:

First, considering man’s inability to always see the truth clearly, differences of opinion about the
outward form of the church are inevitable. Second, even though these differences do not involve
fundamentals of the faith, they are not matters of indifference. Every Christian is obligated to
practice what he believes the Bible teaches. Third, since no church has a final and full grasp of
divine truth, the true church of Christ can never be fully represented by any single ecclesiastical
structure.

Finally, the mere fact of separation does not of itself constitute schism. It is possible to be
divided at many points and still be united in Christ.

Thus the denominational theory of the church looked for Christian unity in some inward religious
experience and allowed diversity in the outward expressions of that personal faith.

This tolerant attitude was not born of doctrinal indifference. The Independent had no intention of
extending Christian unity to all religious professions. The identity of the “one true church” was
restricted to those who shared a common understanding of the core of the Christian faith.

This denominational view of the church found only limited acceptance in England where the Church of
England retained a favored position, even after the Act of Toleration in 1689 recognized the rights
of Presbyterians, Congregationalists, Baptists, and Quakers to worship freely. In the English
colonies of America, however, the denominational theory gained increasing acceptance. It seemed to
be God’s answer for the multiplying faiths in the New World.

Few advocates of the denominational view of the church in the seventeenth century envisioned the
hundreds of Christian groups included under the umbrella today. They had no intention of reducing
the basic beliefs of Christianity to a general feeling of religious sincerity. But they could not
control the future. They simply knew that the traditional bigotry and bloodshed in the name of
Christ was not the way forward.

In the end, then, the denominational form of the church has marked the recent centuries of Christian
history, not because it is ideal, but because it is better than any alternative the years have
offered.

\readinghead{Suggestions for Further Reading}
\begin{readinglist}
\entry Ahlstrom, Sydney E. A Religious History of the American People. New Haven, CT: Yale
  University Press, 1972.
\entry Buschart, W. David. Exploring Protestant Traditions. Downers Grove, IL: IVP Academic, 2006.
\entry *Gaustad, Edwin S. and Leigh E. Schmidt. The Religious History of the America: Revised
  Edition. San Francisco: Harper Collins, 2004.
\entry Hudson, Winthrop S. American Protestantism. Chicago: University of Chicago Press, 1961.
\entry Littell, Franklin H. From State Church to Pluralism. Garden City, NY: Doubleday, 1962.
\entry Mead, Sidney E. The Lively Experiment: The Shaping of Christianity in America. New York:
  Harper \& Row, 1963.
\end{readinglist}
